\appendix
\section{Checklist de conformidade de relato}
\label{ap:checklist}

A Tabela~\ref{tab:checklist} lista os itens de relato recomendados para revisões sistemáticas sob as diretrizes de Kitchenham e Charters \cite{kitchenham2009slr}, indicando, para cada item, onde ele é reportado neste artigo. O objetivo é tornar a conformidade de relato auditável e facilitar a replicação. Os artefatos do protocolo, os logs de triagem, a matriz de extração, a avaliação de qualidade e os scripts de figura e de confiabilidade estão disponíveis no repositório do estudo.

\begin{table}[h!]
\centering
\caption{Checklist de relato Kitchenham e onde cada item é atendido neste artigo.}
\label{tab:checklist}
\small
\begin{tabularx}{\textwidth}{@{}p{0.50\textwidth}Y@{}}
\toprule
\textbf{Item de relato} & \textbf{Onde é atendido} \\
\midrule
Contexto e justificativa da revisão & Introdução \\
Perguntas de pesquisa (RQ1 a RQ5) & Introdução; Método de Revisão \\
Protocolo fixado antes da execução & Método de Revisão; \texttt{protocol.md} \\
Fontes consultadas e strings de busca & Método de Revisão; \texttt{descriptors.md} \\
Critérios de inclusão e exclusão & Método de Revisão (Tabela de critérios) \\
Seleção em duas passadas e contagens & Método de Revisão; logs em \texttt{screening/} \\
Snowballing e análise de saturação & Método de Revisão (Figura de saturação) \\
Confiabilidade entre avaliadores (IRR) & Método de Revisão; \texttt{irr-protocol.md}, \texttt{irr\_kit.py} \\
Esquema e matriz de extração & Método de Revisão; \texttt{extraction-matrix.csv} \\
Avaliação de qualidade / risco de viés & Método de Revisão; Resultados (Tabela de qualidade) \\
Método de síntese & Método de Revisão; Resultados e Síntese \\
Resultados por pergunta de pesquisa & Resultados e Síntese (tabelas tema $\times$ dimensão e tema $\times$ RQ) \\
Ameaças à validade e limitações & Limitações da Revisão \\
Conclusões e trabalho futuro & Agenda de Pesquisa; Conclusão \\
Disponibilidade de dados e artefatos & Repositório do estudo (protocolo, logs, scripts, figuras) \\
\bottomrule
\end{tabularx}
\end{table}
